% ============================================================================
%  Cybersecurity: A Foundations Textbook for AP Cybersecurity : print preamble
%
%  Trim size 7 x 9.25 in
%  twoside,
%  asymmetric inner/outer margins for the binding gutter.
% ============================================================================

\usepackage[T1]{fontenc}
\usepackage[utf8]{inputenc}

% ---------------------------------------------------------------- page setup
\usepackage[
  paperwidth=7in, paperheight=9.25in,
  inner=0.90in, outer=0.70in,
  top=0.80in,   bottom=0.90in,
  headheight=22pt, headsep=10pt,
  footskip=24pt
]{geometry}

% ------------------------------------------------------------------- fonts
%
\usepackage{mathptmx}
\usepackage[scaled=0.90]{helvet}
\usepackage{courier}
\linespread{1.10}        
\usepackage[protrusion=true,expansion=true]{microtype}

% upquote: renders ' and " as straight quotes inside verbatim/listings, so code
% copied out of the PDF is syntactically valid instead of curly-quoted.
\usepackage{upquote}

\usepackage{amsmath,amssymb}
\usepackage{graphicx}

% ------------------------------------------------------- figure/diagram stack
%
%   tikz + pgfplots : network topologies, floor plans, protocol flows, 3D plots
%   bytefield       : packet/header layouts (TCP, IP, DKIM, TLS records)
%   forest          : attack trees and decision trees
%
\usepackage{tikz}
\usepackage{pgfplots}
\pgfplotsset{compat=1.18}
\usetikzlibrary{positioning,arrows.meta,shapes.geometric,calc,fit,backgrounds,patterns,decorations.pathreplacing}
\usepackage{bytefield}
\usepackage{forest}

\usepackage{listings}
\usepackage{xcolor}
\usepackage{tcolorbox}
\usepackage{enumitem}
\usepackage{booktabs}
\usepackage{longtable}
\usepackage{caption}
\usepackage{titlesec}
\usepackage{fancyhdr}
\usepackage{xurl}          % allows long URLs to break across lines
\usepackage{makeidx}       % back-of-book index (imakeidx is not installed here)
\makeindex
\usepackage{hyperref}
\tcbuselibrary{skins,breakable}

% ------------------------------------------------------------------ colours
\definecolor{kickstartermaroon}{HTML}{800000}
\definecolor{kickstartercream}{HTML}{FAF9F6}
\definecolor{codebg}{HTML}{F5F3EE}
\definecolor{rulegray}{HTML}{9A9A9A}

% Extended palette: code syntax highlighting and diagrams.
\definecolor{codecomment}{HTML}{6B5D52}   % warm taupe -- comments, legible, still recessive
\definecolor{codestring}{HTML}{1F4E5F}    % deep steel-teal -- strings; pairs with maroon as the book's second ink colour
\definecolor{codelinenum}{HTML}{8A7F73}   % muted warm gray -- line numbers, present without competing with code
\definecolor{kickstarternavy}{HTML}{1F3864}   % diagram accent 1 -- classic maroon/navy print pairing
\definecolor{kickstarterochre}{HTML}{B5842A}  % diagram accent 2 -- warm highlight, reads clearly on cream
\definecolor{kickstarterforest}{HTML}{2F5D42} % diagram accent 3 -- third hue for multi-part diagrams

% ----------------------------------------------------------------- hyperref
% hidelinks: bookmarks and cross-references stay live in the digital PDF only
\hypersetup{
  hidelinks,
  bookmarksnumbered=true,
  bookmarksopen=true,
  pdftitle={Cybersecurity: A Foundations Textbook for AP Cybersecurity},
  pdfauthor={Mitali Sharma},
  pdfsubject={Cybersecurity},
  pdfcreator={LaTeX}
}

% ------------------------------------------------------- running heads/feet
\pagestyle{fancy}
\fancyhf{}
\renewcommand{\chaptermark}[1]{\markboth{#1}{}}
\renewcommand{\sectionmark}[1]{\markright{#1}}
\fancyhead[LE]{\small\thepage\quad\scshape\nouppercase{\leftmark}}
\fancyhead[RO]{\small\scshape\nouppercase{\rightmark}\quad\thepage}
\renewcommand{\headrulewidth}{0.4pt}
\renewcommand{\headrule}{\hbox to\headwidth{\color{rulegray}\leaders\hrule height \headrulewidth\hfill}}

% Chapter-opening and part pages: centred folio and no running head.
\fancypagestyle{plain}{%
  \fancyhf{}%
  \fancyfoot[C]{\small\thepage}%
  \renewcommand{\headrulewidth}{0pt}%
}

% Truly blank verso pages.
\makeatletter
\def\cleardoublepage{%
  \clearpage
  \if@twoside
    \ifodd\c@page\else
      \hbox{}\thispagestyle{empty}\newpage
      \if@twocolumn\hbox{}\newpage\fi
    \fi
  \fi
}
\makeatother

% ------------------------------------------------------------ sectioning
\titleformat{\part}[display]
  {\normalfont\huge\bfseries\color{kickstartermaroon}}{Part \thepart}{20pt}{\Huge}
\titleformat{\chapter}[display]
  {\normalfont\Large\bfseries\color{kickstartermaroon}}{Chapter \thechapter}{20pt}{\LARGE}
\titlespacing*{\chapter}{0pt}{0pt}{28pt}

% ------------------------------------------------------------- captions
\captionsetup{
  font=small,
  labelfont={bf,sc},
  labelsep=period,
  width=0.92\textwidth
}

% --------------------------------------------------------- code listings
\lstset{
  backgroundcolor=\color{codebg},
  basicstyle=\ttfamily\footnotesize,
  keywordstyle=[1]\color{kickstartermaroon}\bfseries,
  keywordstyle=[2]\color{kickstarterforest},
  morekeywords=[2]{str,int,float,bool,list,dict,set,tuple,bytes,print,len,
    range,enumerate,zip,map,filter,open,type,isinstance,super,sorted,sum,
    min,max,abs,repr,input},
  commentstyle=\itshape\color{codecomment},
  stringstyle=\color{codestring},
  showstringspaces=false,
  upquote=true,              % straight quotes: code stays copy-pasteable
  breaklines=true,
  breakatwhitespace=false,
  % Print a visible continuation arrow when a long line is wrapped, so a
  % reader can tell a wrap from a genuine newline.
  postbreak=\mbox{\textcolor{kickstartermaroon}{$\hookrightarrow$}\space},
  keepspaces=true,
  columns=fullflexible,
  frame=single,
  framerule=0.5pt,
  rulecolor=\color{rulegray},
  numbers=left,
  numberstyle=\tiny\color{codelinenum},
  numbersep=8pt,
  xleftmargin=16pt,
  tabsize=4,
  captionpos=b,
  aboveskip=14pt,
  belowskip=14pt
}

% ------------------------------------------------------------ pedagogy boxes
% Definition box. Every defined term is indexed automatically, so the back-of-
% -book index (see \printindex in main.tex) requires no per-chapter tagging.
\newtcolorbox{definition}[1]{
  colback=kickstartercream, colframe=kickstartermaroon,
  fonttitle=\bfseries, title={Definition: #1}, breakable,
  before upper={\index{#1}}
}

% Misconception / exam-trap box.
\newtcolorbox{misconception}[1][Common Misconception]{
  colback=kickstartercream, colframe=kickstartermaroon!70,
  fonttitle=\bfseries, title={#1}, breakable,
  sharp corners
}

% Worked example box.
\newtcolorbox{example}[1][Worked Example]{
  colback=white, colframe=kickstartermaroon!40,
  fonttitle=\bfseries, title={#1}, breakable
}

\newenvironment{reviewquestions}{\begin{enumerate}[label=\arabic*.]}{\end{enumerate}}

% ------------------------------------------------------------ chapter epigraph
% \chapterquote{quote text}{attribution} -- placed right after \chapter{...}.
% Set in italics, flush right, in the sans body face so it reads as a quote
% rather than as running prose.
\newcommand{\chapterquote}[2]{%
  \begin{flushright}
  \parbox{0.75\textwidth}{\raggedleft\itshape\sffamily\small #1\par
  \normalfont\upshape\footnotesize\MakeUppercase{---~#2}}
  \end{flushright}
  \bigskip
}

%ap align
\newcommand{\chaptertopic}[1]{\noindent{\scshape\color{kickstartermaroon!70}\small AP #1}\par\medskip}

\newtcolorbox{apalignment}[2]{
  colback=kickstartercream, colframe=kickstartermaroon,
  fonttitle=\bfseries\scshape, title={AP Cybersecurity Alignment},
  breakable, sharp corners,
  before upper={\textbf{#1} \,$\cdot$\, #2\par\smallskip}
}

\newcommand{\apobjective}[3]{%
  \item \textbf{\ttfamily \hypertarget{obj:#1}{#1}} \, {\scriptsize\scshape (#2)} --- #3%
}
\newenvironment{apobjectives}{\begin{itemize}[leftmargin=1.2em]}{\end{itemize}
  \par\noindent{\footnotesize Objective codes (e.g., 1.1.A) are this book's own
  numbering, keyed to the Topic number above for cross-reference with the
  review answer key; they are not College Board's numbering. Objective content
  is paraphrased from the \emph{AP Cybersecurity Course Framework} (College
  Board, AP Career Kickstarter).}}

% ------------------------------------------------------------ answer key
% A labelled, numbered list of model answers, printed after reviewquestions.
% Use \answeritem{n}{lo-code}{answer text}.
\newenvironment{answerkey}{%
  \bigskip\noindent\textbf{\scshape Answer Key}\par\smallskip
  \begin{enumerate}[label=\arabic*.]
}{%
  \end{enumerate}
}
% \answerlabel{1.1.A} -- links to the \apobjective with that same code, above.
\newcommand{\answerlabel}[1]{\hfill\hyperlink{obj:#1}{\scriptsize\ttfamily(#1)}\par}

% ------------------------------------------------------- end-of-chapter sources
% \srccite{chapter-key}{n} -- inline citation marker (superscript [n]),
% hyperlinked to the matching \srcitem in that chapter's Sources list.
% \srcitem{chapter-key}{n} -- place at the start of the corresponding \item in
% the Sources enumerate. chapter-key is the chapter's own numeral (e.g. "1.1",
% matching \label{ch:1.1}), so anchor names stay unique across the whole book
% without any extra numbering scheme.
\newcommand{\srccite}[2]{\hyperlink{src:#1:#2}{\textsuperscript{[#2]}}}
\newcommand{\srcitem}[2]{\hypertarget{src:#1:#2}{}}

% ------------------------------------------------------- print housekeeping
\widowpenalty=10000        % never strand a single line at the top of a page
\clubpenalty=10000         % ...or at the bottom
\brokenpenalty=10000       % avoid hyphenating across a page break
\setlength{\emergencystretch}{3em}  % last-resort stretch instead of overfull
\raggedbottom              % better than stretching around code/box floats
